\section{本章小结}

本章把实数和欧氏空间上的长度、体积与积分推广到一般底集，得到四条主要认识。

第一，一般底集上必须另外给定可测范围、赋权和（在需要几何结论时）拓扑结构。
$\sigma$代数规定哪些集合可以测量，生成族使可测性和测度唯一性只需在简单集合上检验；测度规定这些集合的非负权重，有限、概率和$\sigma$有限是不同强度的质量条件；Borel结构、Polish空间和支持把可测集合与开闭邻域联系起来。
分层积分以及单调收敛、Fatou和支配收敛的机制仍沿用微积分，乘积结构、Tonelli/Fubini和柱集接口则把积分推广到联合对象。

第二，固定同一空间时，零集关系说明两项赋权是否相容。
绝对连续、等价和相互奇异分别刻画零集的单向蕴含、双向共享和质量范围分离；Radon–Nikodym密度恢复绝对连续部分的权重，Lebesgue分解把不能用密度表示的奇异部分另外拆出；有符号测度的Hahn–Jordan分解和总变差累计正负变化，不允许它们抵消。
全变差、Hellinger、KL和Wasserstein从不同机制量化概率差异：前三者在固定位置比较事件或密度，Wasserstein还要使用共同距离；积分概率度量把这些指标统一成“某类测试函数的最大积分差”，而不是再添加一个互斥的指标类别。

第三，跨空间时，可测映射为每个输入指定一个目标点，测度核为每个输入指定一项目标测度。
在输入测度有限、核有统一有限质量界的范围内，核积分给出输出测度并保持观测传递，核复合扩展了映射复合；联合测度同时记录输入和输出，乘积和图测度都是其特例，沿第一边缘在标准Borel条件下可反向分解出归一化核。
核的总质量可以随输入变化，归一化条件才保证质量保持。
较粗可测结构上的条件积分建立了通用的积分保持关系；条件期望的代数法则、Bayes公式和随机过程解释由后文概率论继续展开。

第四，建立共同空间后比较的是输出权重，给定跨空间成本后优化的是配对代价。
同一归一化测度核收缩全变差，但这一结论不能无条件推广到两个不同的传递规则；表示后的零差异只说明该表示看不到某些输入区别，不能恢复丢掉的信息。
跨空间运输要求先固定可测的成本函数；Kantorovich计划允许拆分质量，Monge映射不允许，最优计划达到还要依赖一致紧性和成本下半连续。
比较之前应先固定对应关系、允许的观测和成本，不能只凭两侧内部距离就直接计算跨空间距离。

第五，序列与曲线的极限依赖所用标准。
逐集、全变差、弱收敛和Wasserstein收敛依次放宽或加入不同结构；弱收敛由有界连续测试定义，Portmanteau界控制边界零质量事件，连续推前保持弱极限，下半连续成本只给积分下界，一致紧性经Prokhorov定理保证弱子列，Wasserstein收敛还要额外控制距离尾部。
后文的随机过程、条件模型、经验测度和策略分析都可以调用这些通用操作，但它们的概率性质、统计保证和生成机制仍需各自的条件。
